\chapter{Discussion and Conclusions}\label{discussion}
\section{Introduction}
This chapter discusses the findings presented in Chapter \ref{results} in relation to the research objectives, literature review, and design principles developed through the Design Science Research process. The chapter evaluates the extent to which the artefact addressed the identified problem, reflects upon the strengths and limitations of the solution, and considers opportunities for future development.

\section{Discussion of Research Question 1}
\textbf{Where do delays occur between MDT outcome and treatment?}

The primary objective of this study was to improve visibility of delays occurring between MDT outcome and treatment initiation. Analysis of the integrated pathway demonstrated that delays were not distributed evenly across the pathway. Rather, two stages accounted for the majority of elapsed time: triage-to-clinic attendance (8 days) and CT simulation-to-treatment initiation (16 days). 

The triage-to-clinic interval represents a period during which referrals are reviewed, prioritised, booked, and ultimately seen by an oncologist. Although the observed interval remained well within national cancer waiting time standards, it nevertheless represented a meaningful proportion of the overall pathway duration. This finding suggests that clinic access capacity remains an important determinant of pathway flow and warrants continued operational monitoring. 

The CT simulation-to-treatment interval represented the largest component of the pathway. However, unlike other pathway stages, this interval does not correspond to a single operational activity. Instead, it encompasses multiple planning processes including target contouring, treatment planning, peer review, quality assurance activities, and treatment booking. Consequently, the current analysis identifies where time accumulates but cannot yet determine which individual planning activities contribute most significantly to delay. 

These findings support the argument presented within the literature that aggregated treatment interval metrics provide limited understanding of operational bottlenecks. By decomposing the pathway into canonical events, the artefact provided substantially greater visibility than traditional compliance reporting.

\section{Discussion of Research Question 2}
\textbf{What operational factors are associated with delay?}

A secondary objective of this study was to identify operational factors associated with delay. Although the repository was not designed as a predictive modelling platform, several factors emerged during analysis. 

First, increasing future clinic waits were observed despite relatively short historic referral-to-clinic intervals. Patients holding future appointments were waiting substantially longer than patients who had already completed the clinic stage. Although future booked waits were considerably longer than the referral-to-clinic interval, these waits largely reflected planned follow-up and review activity rather than new-patient access. The finding therefore provides greater insight into future clinic workload than referral-processing performance. Further investigation, incorporating clinic templates and appointment-type classifications, would be required before conclusions could be drawn regarding clinic-capacity constraints affecting new referrals.

Second, data-quality issues emerged as an operational factor in their own right. A substantial number of clinic appointments lacked recorded outcomes, while duplicate referral records and inconsistent pathway representations required extensive harmonisation during ETL processing. These findings indicate that administrative recording practices can influence both operational visibility and service-management effectiveness. 

Third, radiotherapy demand was observed to operate close to departmental planning thresholds. Although utilisation remained within operationally manageable levels, average projected utilisation exceeded the internal planning threshold of 80\%. This finding highlights the importance of balancing efficiency with sufficient flexibility to accommodate urgent treatments and unplanned events.

\section{Artefact Evaluation}
The artefact successfully achieved its primary objective of providing integrated operational visibility across previously disconnected stages of the post-MDT oncology pathway. 

From a technical perspective, the repository successfully integrated data from multiple heterogeneous operational systems using a common ETL framework. Automated refresh processes reduced manual data processing requirements while preserving data provenance and auditability. 

From an operational perspective, the dashboard suite provided users with visibility of referral activity, clinic performance, radiotherapy capacity, and end-to-end pathway progression. The separation of functionality across multiple dashboards appeared to support task-focused operational workflows without overwhelming users with excessive information. 

Importantly, the development process demonstrated that substantial value can be obtained from relatively lightweight integration architectures. Rather than replacing existing systems, the artefact created a unified analytical layer across existing operational data sources.

\section{Evaluation of Design Principles}
The Design Principles developed in Chapter \ref{method} were intended to represent transferable knowledge extending beyond the specific implementation used in this study. 

DP1 and DP2 received particularly strong support during evaluation. Risk-based patient categorisation and capacity-demand visualisation enabled operational users to identify not only where delays existed but also whether sufficient operational capacity existed to mitigate those risks. 

DP3 was supported through the use of canonical pathway events aligned to published cancer waiting time standards. The pathway model enabled consistent interval calculation across multiple source systems and reduced ambiguity associated with locally defined metrics. 

DP4 was partially supported. While daily automated refreshes were achieved for the majority of datasets, dependency upon external operational reports limited refresh frequency for some clinic data extracts. Nevertheless, the data remained sufficiently current to support operational decision-making. 

DP5 was supported through the separation of staging, transformation, and reporting layers, allowing all reported metrics to be traced back to source records. 

DP6 was supported through iterative stakeholder engagement and the creation of task-oriented dashboards aligned with referral management, radiotherapy service management, and pathway performance monitoring respectively. 

Collectively, the evaluation suggests that all six design principles remain valid and may have applicability beyond the specific organisational setting studied.

\section{Reflections and Lessons Learned}
Several important lessons emerged during the development process. 

The most significant lesson related to data quality. Initial assumptions suggested that the primary challenge would involve technical integration of multiple systems. In practice, the greater challenge involved understanding operational processes and interpreting inconsistently recorded data. 

The project also highlighted the value of Design Science Research for healthcare informatics projects. The iterative nature of the methodology enabled technical development to remain aligned with operational requirements while simultaneously generating transferable design knowledge. 

Finally, the study reinforced the importance of stakeholder engagement. Several dashboard features and pathway metrics evolved considerably following operational feedback. This iterative refinement improved both the utility and relevance of the final artefact.

\section{Limitations}

Several limitations should be considered when interpreting the findings. 

First, the study was conducted within a single NHS organisation and therefore reflects local processes, systems, and operational practices. 

Second, the repository relies upon the quality and completeness of source-system data. Although significant harmonisation and validation activities were undertaken, residual inaccuracies within source datasets may still influence results. 

Third, the CT simulation-to-treatment interval represents an aggregation of multiple planning and operational activities. While the repository successfully identifies this stage as a significant contributor to overall pathway duration, the current data model cannot determine precisely how time is distributed between individual planning processes. 

Finally, evaluation focused primarily on operational utility and usability rather than direct patient outcomes. Consequently, although the artefact provides improved operational visibility, the extent to which it contributes to improved clinical outcomes remains outside the scope of the present study.

\section{Future Work}
\subsection{Clinic Capacity Modelling}
Future development should incorporate clinic-capacity modelling and appointment-type classification to better understand demand distribution across new-patient, follow-up, and review clinics. This would provide greater visibility of capacity utilisation and allow more detailed investigation of factors influencing access to first oncology appointments.

\subsection{Digital Referral Forms}

The project identified several data-quality issues attributable to manually maintained processes and retrospective transcription of referral information. Future work should investigate the implementation of structured digital referral forms containing standardised mandatory fields and validation rules. Such an approach may improve data quality while reducing ETL complexity.

\subsection{CT-to-Treatment Decomposition}
The current pathway model treats CT simulation-to-treatment as a single interval. Future repository development should introduce additional canonical events representing contouring, treatment planning, peer review, quality assurance, and treatment booking activities. This would provide a more granular understanding of where delays occur within the treatment-planning process.

\subsection{SACT Integration}
A key future enhancement involves integration of Systemic Anti-Cancer Therapy (SACT) data following completion of the planned Cancer Network-wide system upgrade. Incorporating SACT treatment events would enable end-to-end analysis of all major treatment modalities and provide a more comprehensive representation of post-MDT oncology pathways.

\subsection{Predictive Analytics}
The repository architecture also provides a foundation for future predictive analytics. Potential applications include forecasting breach risk, modelling future demand, and identifying patients likely to experience pathway delays before breaches occur.

\section{Conclusion}
This study has demonstrated that meaningful operational insight can be generated through the integration of existing oncology datasets within a lightweight Integrated Data Repository. Using a Design Science Research methodology, the project successfully developed and evaluated an artefact capable of improving visibility of post-MDT pathway performance. 

The repository identified that delays are concentrated primarily within the triage-to-clinic and CT simulation-to-treatment stages of the pathway. Furthermore, the project demonstrated that significant operational value can be achieved through relatively modest data integration and analytics capabilities without requiring replacement of existing clinical systems. 

Beyond the artefact itself, the study contributes a set of evaluated design principles for operational oncology analytics that may be transferable to other healthcare organisations facing similar challenges relating to fragmented data and pathway visibility. 

The findings therefore contribute both practical value for operational cancer services and academic value through the generation of reusable design knowledge for healthcare informatics.